本研究では、欠陥予測を「その変更に欠陥が含まれるかどうか」という二値分類問題として扱った。しかし、ソフトウェア開発においては、欠陥の混入と修正は必ずしも1対1ではなく、多対多の対応関係を持つ。本節では、この多対多関係が予測モデルに与える影響と、その解決策について考察する。

\paragraph{多対多関係の実態と課題}
実際のプロジェクトでは、1つの大きな変更が複数の独立した欠陥を誘発するケース（1:N）や、複数のコミットにまたがる変更が積み重なって初めて1つの欠陥が顕在化するケース（N:1）が発生する。現在の二値分類による欠陥予測では、これらの依存関係を「欠陥の有無」という単一のラベルに集約することになり、各コミットが持つ欠陥混入への寄与度を詳細に評価できないという課題がある。

\paragraph{回帰モデルによる寄与度の予測}
この多対多関係を適切に扱うための解決策として、二値分類から回帰モデルへの移行が有効である。回帰モデルを導入することで、特定のコミットが「いくつの欠陥に関与したか」あるいは「欠陥全体に対してどの程度の責任（寄与度）を持つか」を連続値として予測可能になる。これにより、3.4節で定義した密度 $D_{i}$ において、「予測された欠陥の影響度」を $\hat{y}_{i}$ として活用でき、レビューの優先順位付けをより正確に行える。

\paragraph{SZZアルゴリズムのノイズと相対的リスクの識別}
回帰モデルの導入にあたっては、正解ラベルの生成に用いるSZZアルゴリズムの精度が課題となる。SZZは欠陥修正の見落としや誤検出を含むため、欠陥混入への寄与度を厳密な連続値として定義しようとすると、学習データに大きなノイズが生じる懸念がある。
しかし、開発現場の意思決定支援においては、欠陥混入数の予測よりも、リソース配分のための「相対的な欠陥混入リスクの差異」を正しく捉えることがより重要である。したがって、欠陥混入数そのものではなく、プロジェクト内での相対的な影響度を予測することで、SZZの不完全性を許容しつつ、多対多の環境下でも高リスクな変更を効果的に識別できると考えられる。

\paragraph{説明可能性を考慮したアルゴリズム選定}
多対多の関係性をモデル化する場合であっても、本研究が重視している「予測の根拠」の提示は不可欠である。しかし、分類と回帰の両方に対応している機械学習アルゴリズムは多い。一例として、ランダムフォレストを用いた回帰モデルでは、特徴量重要度や決定木の可視化を通じた分析が可能である。これにより、複雑な依存関係においてどのメトリクスが欠陥混入の主な原因となっているかを、開発者が直感的に理解できる形でフィードバックすることが可能となる。